\section{Decision problem and roles}
\label{sec:framework}

An instance is a latent site $\theta\sim P_\Theta$ defining an objective
$f_\theta:\mathcal X\rightarrow\mathbb R$. In round $r$, the agent observes a timestamped
history $H_r$, chooses a finite design $D_r$, and reports a recommendation $\hat x_r$. A
design contains factor settings and an assignment to available physical units
$\mathcal U_r$; feasibility therefore depends on the facility, not only on whether
$x\in\mathcal X$.

The \emph{design policy} $\pi_D$ maps visible history and remaining resources to a design,
\begin{equation}
 D_r=\pi_D(H_r,\mathcal U_r)\in\mathcal D.
 \label{eq:design}
\end{equation}
It controls what evidence can exist. The \emph{history reader} $\rho$ maps the same visible
history to beliefs and a recommendation,
\begin{equation}
 \hat x_r=\rho(H_r)\in\mathcal X.
 \label{eq:reader}
\end{equation}
Equations~\ref{eq:design} and~\ref{eq:reader} make the separation explicit: the first role controls what evidence can exist, while the second controls how acquired evidence is used. These are analytical roles rather than required software modules; the same LLM produces both outputs in a closed-loop episode.

Submitting $D_r$ occupies its units for a fixed duration $\Delta$. Measurements enter
$H_t$ only when their timestamps are reached. The benchmark permits inspection during a
cycle but does not permit treatment modification, early release, or refunds. Terminal
outcomes follow
\begin{equation}
 y_u=f_\theta(x_u;\phi_u)+\varepsilon_u,
 \qquad \varepsilon\sim\mathcal N(0,\Sigma(D_r)),
 \label{eq:observation}
\end{equation}
Equation~\ref{eq:observation} makes the observation hierarchy explicit: per-unit biological variation $\phi_u$ and chamber, loop, batch, and measurement noise make the covariance structured and treatment dependent. A poor design can therefore fail by
choosing uninformative treatments, confounding them with blocks, or allocating too little
replication.

The \emph{evaluator} is outside the episode policy. It receives designs and recommendations
and may access $\theta$, $f_\theta$, $P_\Theta$, and the forward model. These objects never
enter $H_t$, and evaluator outputs are not fed back to the agent. A design aid may change
$\pi_D$; an inference aid may change $\rho$; neither is an evaluator.

The four regime properties follow directly. Feedback is \emph{slow} because finite
parallel capacity and latency bound evidence per calendar time. It is \emph{noisy} because
one unit does not settle a comparison. It is \emph{costly} because experiments earn the
same economic objective they seek to optimise. It is \emph{irreversible} because committed
capacity and treatments cannot be recovered after later evidence arrives.
